% Generated from lessons/0007-3-fundamental-theorem-calculus.html; edit the HTML source, then rerun the converter.
\section{微积分基本定理}
\lessonmark{微积分基本定理}

本节把“积分是极限和”与“导数是变化率”接起来。考试中最常见的动作是：变上限定积分求导、Newton-Leibniz 计算定积分、换元、分部积分，以及把求和极限看成 Riemann 和。

\subsection*{本节主线}

\begin{conceptbox}
\textbf{积分产生原函数}
\(F(x)=\int_a^x f(t)\,dt\) 在 \(f\) 连续时满足 \(F'(x)=f(x)\)。
\end{conceptbox}

\begin{conceptbox}
\textbf{原函数计算积分}
若 \(\Phi'(x)=f(x)\)，则 \(\int_a^b f(x)\,dx=\Phi(b)-\Phi(a)\)。
\end{conceptbox}

\begin{conceptbox}
\textbf{计算技术}
换元和分部积分，本质是链式法则与乘积求导法则的积分版。
\end{conceptbox}

\begin{notebox}
这节最重要的感觉：变上限定积分求导时，“上限动到哪里，就把被积函数代到哪里”，再乘以上限的导数；下限同理但带负号。
\end{notebox}

\subsection*{一、变上限定积分}

\subsubsection*{基本定理第一部分}

若 \(f\) 在 \([a,b]\) 上连续，定义

\[
      F(x)=\int_a^x f(t)\,dt,\qquad x\in[a,b],
    \]

则 \(F\) 在 \((a,b)\) 内可导，且

\[
      F'(x)=f(x).
    \]

\textbf{证明思路：}

\[
      \frac{F(x+h)-F(x)}{h}
      =
      \frac1h\int_x^{x+h}f(t)\,dt.
    \]

右边是小区间 \([x,x+h]\) 上的平均值；当 \(h\to0\) 时，由 \(f\) 的连续性，平均值趋于 \(f(x)\)。

\subsubsection*{上下限都是函数}

若 \(f\) 连续，\(\alpha,\beta\) 可导，则

\[
      \frac{d}{dx}\int_{\alpha(x)}^{\beta(x)}f(t)\,dt
      =
      f(\beta(x))\beta'(x)-f(\alpha(x))\alpha'(x).
    \]

这是考试最常用公式。上限贡献正号，下限贡献负号。

\subsection*{二、Newton-Leibniz 公式}

若 \(f\) 在 \([a,b]\) 上连续，\(\Phi\) 是 \(f\) 的一个原函数，即 \(\Phi'(x)=f(x)\)，则

\[
      \int_a^b f(x)\,dx=\Phi(b)-\Phi(a).
    \]

\textbf{证明思路：}令 \(F(x)=\int_a^x f(t)\,dt\)。由第一部分 \(F'(x)=f(x)\)，所以 \(F\) 与 \(\Phi\) 只差一个常数。代入端点即可。

\begin{notebox}
计算定积分时，原函数只是工具；最终答案是端点差，不要忘记代上下限。
\end{notebox}

\subsection*{三、换元积分法}

若 \(f\) 连续，\(\varphi\) 有连续导数，则

\[
      \int_\alpha^\beta f(\varphi(t))\varphi'(t)\,dt
      =
      \int_{\varphi(\alpha)}^{\varphi(\beta)} f(u)\,du.
    \]

\textbf{证明思路：}若 \(F'(u)=f(u)\)，则 \([F(\varphi(t))]'=f(\varphi(t))\varphi'(t)\)。这就是链式法则。

\begin{center}
\small
\begin{tabularx}{\linewidth}{|>{\raggedright\arraybackslash}X|>{\raggedright\arraybackslash}X|}
\hline
\textbf{常见换元信号} & \textbf{尝试} \\ \hline
根式 \(a^2-x^2\) & 三角换元 \(x=a\sin t\) 或 \(x=a\cos t\) \\ \hline
\(\ln x\) & 令 \(u=\ln x\)，常配合 \(dx/x\) \\ \hline
对称区间 & 令 \(x=a+b-t\)，把区间反过来 \\ \hline
Riemann 和 & 识别 \(\frac1n\sum f(k/n)\) \\ \hline
\end{tabularx}
\end{center}

\subsection*{四、分部积分法}

若 \(u,v\) 有连续导数，则

\[
      \int_a^b u(x)v'(x)\,dx
      =
      \bigl[u(x)v(x)\bigr]_a^b-\int_a^b u'(x)v(x)\,dx.
    \]

\textbf{证明思路：}由 \((uv)'=u'v+uv'\)，两边积分。

\begin{notebox}
看到 \(x\sin x\)、\(\ln x\)、\(\arctan x\)、\(x^m e^x\) 这类“一个代数因子乘一个容易积分/求导的因子”，优先想到分部积分。
\end{notebox}

\subsection*{五、对称性、周期性与 Riemann 和}

\subsubsection*{对称替换}

\[
      \int_a^b f(x)\,dx=\int_a^b f(a+b-x)\,dx.
    \]

这条来自换元 \(u=a+b-x\)。它常用于把含 \(x\) 的积分和含 \(a+b-x\) 的积分相加。

\[
      \int_0^\pi x f(\sin x)\,dx
      =
      \frac{\pi}{2}\int_0^\pi f(\sin x)\,dx.
    \]

\subsubsection*{三角正交性}

在 \([-\pi,\pi]\) 上，三角函数族具有正交性：

\[
      \int_{-\pi}^{\pi}\sin mx\cos nx\,dx=0,
    \]

\[
      \int_{-\pi}^{\pi}\sin mx\sin nx\,dx=
      \begin{cases}
      0,&m\ne n,\\
      \pi,&m=n,
      \end{cases}
      \qquad
      \int_{-\pi}^{\pi}\cos mx\cos nx\,dx=
      \begin{cases}
      0,&m\ne n,\\
      \pi,&m=n\ne0,\\
      2\pi,&m=n=0.
      \end{cases}
    \]

\subsubsection*{Riemann 和极限}

若形如

\[
      \frac1n\sum_{k=1}^n f\!\left(\frac{k}{n}\right),
    \]

通常对应

\[
      \int_0^1 f(x)\,dx.
    \]

关键是识别“采样点” \(\frac{k}{n}\) 和“小宽度” \(\frac1n\)。

\subsection*{六、课后题训练}

下面按教材第 264-267 页习题 1-26 整理。题面保留，提示只给入口；写作业时建议先独立算一遍，再展开提示核对方向。

\begin{exercisebox}
\textbf{习题 1 \badge{变上限求导}}

设函数 \(f(x)\) 连续，求下列函数 \(F(x)\) 的导数：

\[
      \begin{aligned}
      &(1)\ F(x)=\int_x^b f(t)\,dt;\\
      &(2)\ F(x)=\int_a^{\ln x} f(t)\,dt;\\
      &(3)\ F(x)=\int_a^{\int_0^x\sin^2s\,ds}\frac{1}{1+t^2}\,dt.
      \end{aligned}
      \]

\begin{hintbox}
上限贡献正号，下限贡献负号；第 (3) 小题是两层链式法则。
\end{hintbox}
\end{exercisebox}

\begin{exercisebox}
\textbf{习题 2 \badge{极限与 L'Hospital}}

求下列极限：

\[
      \begin{aligned}
      &(1)\ \lim_{x\to0}\frac{\int_0^x\cos t^2\,dt}{x};\\
      &(2)\ \lim_{x\to0}\frac{x^2}{\int_{\cos x}^{1}e^{-1/w^2}\,dw};\\
      &(3)\ \lim_{x\to+\infty}\frac{\int_0^x(\arctan v)^2\,dv}{\sqrt{1+x^2}};\\
      &(4)\ \lim_{x\to+\infty}\frac{\left(\int_0^x e^{u^2}\,du\right)^2}{\int_0^x e^{2u^2}\,du}.
      \end{aligned}
      \]

\begin{hintbox}
分子或分母含变上限定积分时，求导会把被积函数“吐出来”。第 (4) 小题要小心链式法则。
\end{hintbox}
\end{exercisebox}

\begin{exercisebox}
\textbf{习题 3 \badge{单调性}}

设 \(f(x)\) 是 \([0,+\infty)\) 上的连续函数且恒有 \(f(x)>0\)，证明

\[
        g(x)=\frac{\int_0^x t f(t)\,dt}{\int_0^x f(t)\,dt}
      \]

是定义在 \([0,+\infty)\) 上的单调增加函数。

\begin{hintbox}
把 \(g(x)\) 看成权重为 \(f(t)\) 的平均值；也可直接对商求导，再把分子整理成双重积分或“平均值差”。
\end{hintbox}
\end{exercisebox}

\begin{exercisebox}
\textbf{习题 4 \badge{极值}}

求函数

\[
        f(x)=\int_0^x(t-1)(t-2)^2\,dt
      \]

的极值。

\begin{hintbox}
由微积分基本定理，先求 \(f'(x)\)，再研究临界点附近符号。
\end{hintbox}
\end{exercisebox}

\begin{exercisebox}
\textbf{习题 5 \badge{中值定理求极限}}

利用中值定理求下列极限：

\[
      (1)\ \lim_{n\to\infty}\int_0^1\frac{x^n}{1+x}\,dx,
      \qquad
      (2)\ \lim_{n\to\infty}\int_n^{n+p}\frac{\sin x}{x}\,dx\quad(p\in\mathbb N^*).
      \]

\begin{hintbox}
第 (1) 题可用夹逼；第 (2) 题区间长度固定但 \(1/x\to0\)。
\end{hintbox}
\end{exercisebox}

\begin{exercisebox}
\textbf{习题 6 \badge{定积分计算}}

求下列定积分：

\[
      \begin{aligned}
      &(1)\ \int_0^1 x^2(2-x^2)^2\,dx;\\
      &(2)\ \int_1^2\frac{(x-1)(x^2-x+1)}{2x^2}\,dx;\\
      &(3)\ \int_0^2(2^x+3^x)^2\,dx;\\
      &(4)\ \int_0^{1/2}x(1-4x^2)^{10}\,dx;\\
      &(5)\ \int_{-1}^{1}\frac{(x+1)\,dx}{(x^2+2x+5)^2};\\
      &(6)\ \int_0^1\arcsin x\,dx;\\
      &(7)\ \int_{-\pi/4}^{\pi/4}\frac{x}{\cos^2x}\,dx;\\
      &(8)\ \int_0^{\pi/4}x\tan^2x\,dx;\\
      &(9)\ \int_0^{\pi/2}e^x\sin^2x\,dx;\\
      &(10)\ \int_1^e\sin(\ln x)\,dx;\\
      &(11)\ \int_0^1 x^2\arctan x\,dx;\\
      &(12)\ \int_1^{e+1}x^2\ln(x-1)\,dx;\\
      &(13)\ \int_0^{\sqrt{\ln2}}x^3e^{-x^2}\,dx;\\
      &(14)\ \int_0^1 e^{2\sqrt{x+1}}\,dx;\\
      &(15)\ \int_0^1\frac{dx}{\sqrt{1+e^{2x}}};\\
      &(16)\ \int_{-1/2}^{1/2}\frac{dx}{\sqrt{(1-x^2)^3}};\\
      &(17)\ \int_0^1\left(\frac{x-1}{x+1}\right)^4dx;\\
      &(18)\ \int_0^1\frac{x^2+1}{x^4+1}\,dx;\\
      &(19)\ \int_1^{\sqrt2}\frac{dx}{x\sqrt{1+x^2}};\\
      &(20)\ \int_0^1x\sqrt{\frac{x}{2-x}}\,dx.
      \end{aligned}
      \]

\begin{hintbox}
先分类：多项式展开、换元、分部积分、奇偶性/对称性、三角换元。不要一上来硬算。
\end{hintbox}
\end{exercisebox}

\begin{exercisebox}
\textbf{习题 7 \badge{Riemann 和}}

求下列极限：

\[
      \begin{aligned}
      &(1)\ \lim_{n\to\infty}\left(\frac1{n^2}+\frac2{n^2}+\frac3{n^2}+\cdots+\frac{n-1}{n^2}\right);\\
      &(2)\ \lim_{n\to\infty}\frac{1^p+2^p+\cdots+n^p}{n^{p+1}}\quad(p>0);\\
      &(3)\ \lim_{n\to\infty}\frac1n\left(\sin\frac{\pi}{n}+\sin\frac{2\pi}{n}+\cdots+\sin\frac{(n-1)\pi}{n}\right).
      \end{aligned}
      \]

\begin{hintbox}
把每一项写成 \(f(k/n)\cdot(1/n)\)。
\end{hintbox}
\end{exercisebox}

\begin{exercisebox}
\textbf{习题 8 \badge{参数积分}}

求下列定积分：

\[
      \begin{aligned}
      &(1)\ \int_0^\pi \cos^n x\,dx;\\
      &(2)\ \int_{-\pi}^{\pi}\sin^n x\,dx;\\
      &(3)\ \int_0^a(a^2-x^2)^n\,dx;\\
      &(4)\ \int_0^{1/2}x^2(1-4x^2)^{10}\,dx;\\
      &(5)\ \int_0^1 x^n\ln^m x\,dx;\\
      &(6)\ \int_1^e x\ln^n x\,dx.
      \end{aligned}
      \]

\begin{hintbox}
第 (1)(2) 题分 \(n\) 奇偶；第 (5) 题可反复分部积分或令 \(x=e^{-t}\)。
\end{hintbox}
\end{exercisebox}

\begin{exercisebox}
\textbf{习题 9 \badge{对称性}}

设 \(f(x)\) 在 \([0,1]\) 上连续，证明：

\[
      (1)\ \int_0^{\pi/2}f(\cos x)\,dx=\int_0^{\pi/2}f(\sin x)\,dx,
      \qquad
      (2)\ \int_0^\pi x f(\sin x)\,dx=\frac{\pi}{2}\int_0^\pi f(\sin x)\,dx.
      \]

\begin{hintbox}
第 (1) 题令 \(x=\pi/2-t\)；第 (2) 题令 \(x=\pi-t\)，再与原式相加。
\end{hintbox}
\end{exercisebox}

\begin{exercisebox}
\textbf{习题 10 \badge{对称性计算}}

利用上题结果计算：

\[
      (1)\ \int_0^\pi x\sin^4x\,dx,\quad
      (2)\ \int_0^\pi \frac{x\sin x}{1+\cos^2x}\,dx,\quad
      (3)\ \int_0^\pi \frac{x}{1+\sin^2x}\,dx.
      \]

\begin{hintbox}
先把 \(x\) 消掉，变成 \(\frac{\pi}{2}\) 倍普通三角积分。
\end{hintbox}
\end{exercisebox}

\begin{exercisebox}
\textbf{习题 11 \badge{分段函数}}

求下列定积分：

\[
      (1)\ \int_0^6 x^2[x]\,dx,\quad
      (2)\ \int_0^2\operatorname{sgn}(x-x^3)\,dx,\quad
      (3)\ \int_0^1 x|x-a|\,dx,\quad
      (4)\ \int_0^2[e^x]\,dx.
      \]

\begin{hintbox}
找分段点：整数点、零点、绝对值折点、取整函数跳跃点。
\end{hintbox}
\end{exercisebox}

\begin{exercisebox}
\textbf{习题 12 \badge{对称函数}}

设 \(f(x)\) 在 \([a,b]\) 上可积且关于 \(x=T\) 对称，这里 \(a<T<b\)。证明

\[
        \int_a^b f(x)\,dx=\int_a^{2T-b} f(x)\,dx+2\int_T^b f(x)\,dx,
      \]

并给出几何解释。

\begin{hintbox}
把 \([2T-b,b]\) 关于 \(T\) 配对，左、右面积相等。
\end{hintbox}
\end{exercisebox}

\begin{exercisebox}
\textbf{习题 13 \badge{分段与平移}}

设

\[
      f(x)=
      \begin{cases}
      xe^{-x^2},&x\ge0,\\
      \frac1{1+e^x},&x<0,
      \end{cases}
      \]

计算 \(I=\int_1^4 f(x-2)\,dx\)。

\begin{hintbox}
令 \(u=x-2\)，区间变为 \([-1,2]\)，再在 \(u=0\) 分段。
\end{hintbox}
\end{exercisebox}

\begin{exercisebox}
\textbf{习题 14 \badge{含参积分求导}}

设函数

\[
      f(x)=\frac12\int_0^x (x-t)^2g(t)\,dt,
      \]

其中 \(g(x)\) 在 \((-\infty,+\infty)\) 上连续，且 \(g(1)=5,\int_0^1g(t)\,dt=2\)。证明

\[
      f'(x)=x\int_0^xg(t)\,dt-\int_0^x tg(t)\,dt,
      \]

并计算 \(f''(1)\) 和 \(f'''(1)\)。

\begin{hintbox}
可以先展开 \((x-t)^2\)，也可以用 Leibniz 公式。注意上限项中 \((x-x)^2=0\)。
\end{hintbox}
\end{exercisebox}

\begin{exercisebox}
\textbf{习题 15 \badge{积分方程}}

设 \((0,+\infty)\) 上连续函数 \(f(x)\) 满足

\[
      f(x)=\ln x-\int_1^x f(t)\,dt,
      \]

求 \(\int_1^e f(x)\,dx\)。

\begin{hintbox}
令 \(F(x)=\int_1^x f(t)\,dt\)，把题目变成关于 \(F\) 的一阶微分方程。
\end{hintbox}
\end{exercisebox}

\begin{exercisebox}
\textbf{习题 16 \badge{积分方程换元}}

设函数 \(f(x)\) 连续，且

\[
      \int_0^1tf(2x-t)\,dt=\frac12\arctan(x^2),\qquad f(1)=1.
      \]

求 \(\int_1^2 f(x)\,dx\)。

\begin{hintbox}
对 \(x\) 求导或令 \(u=2x-t\)，把未知积分区间显露出来。
\end{hintbox}
\end{exercisebox}

\begin{exercisebox}
\textbf{习题 17 \badge{周期绝对值}}

求

\[
      \int_0^{n\pi}x|\sin x|\,dx,
      \]

其中 \(n\) 为正整数。

\begin{hintbox}
按每个 \([k\pi,(k+1)\pi]\) 分段，令 \(x=k\pi+t\)。
\end{hintbox}
\end{exercisebox}

\begin{exercisebox}
\textbf{习题 18 \badge{平均增长率}}

设函数

\[
      S(x)=\int_0^x|\cos t|\,dt,
      \]

求 \(\lim_{x\to+\infty}S(x)/x\)。

\begin{hintbox}
\(|\cos t|\) 的周期为 \(\pi\)。长期平均值等于一个周期积分除以周期长度。
\end{hintbox}
\end{exercisebox}

\begin{exercisebox}
\textbf{习题 19 \badge{尺度不变}}

设 \(f(x)\) 在 \((0,+\infty)\) 上连续，且对于任意 \(a>0\) 有

\[
      g(x)=\int_x^{ax}f(t)\,dt=\text{常数},\qquad x\in(0,+\infty).
      \]

证明 \(f(x)=c/x\)，其中 \(c\) 为常数。

\begin{hintbox}
对 \(x\) 求导：上限 \(ax\) 与下限 \(x\) 都会产生项。
\end{hintbox}
\end{exercisebox}

\begin{exercisebox}
\textbf{习题 20 \badge{对称换元}}

设 \(f(x)\) 在 \((0,+\infty)\) 上连续，证明：

\[
      \int_1^4 f\!\left(\frac{x}{2}+\frac{2}{x}\right)\frac{\ln x}{x}\,dx
      =
      (\ln2)\int_1^4 f\!\left(\frac{x}{2}+\frac{2}{x}\right)\frac{1}{x}\,dx.
      \]

\begin{hintbox}
令 \(x=4/t\)，观察 \(\frac{x}{2}+\frac{2}{x}\) 不变，而 \(\ln x\) 与 \(\ln 4-\ln t\) 互补。
\end{hintbox}
\end{exercisebox}

\begin{exercisebox}
\textbf{习题 21 \badge{估计不等式}}

设 \(f'(x)\) 在 \([a,b]\) 上连续，证明：

\[
      \max_{a\le x\le b}|f(x)|
      \le
      \left|\frac1{b-a}\int_a^b f(x)\,dx\right|
      +\int_a^b|f'(x)|\,dx.
      \]

\begin{hintbox}
先把 \(f(x)\) 写成区间平均值加一个偏差；偏差用 \(f'\) 控制。
\end{hintbox}
\end{exercisebox}

\begin{exercisebox}
\textbf{习题 22 \badge{二重积分影子}}

设 \(f(x)\) 在 \((-\infty,+\infty)\) 上连续，证明：

\[
      \int_0^x f(u)(x-u)\,du
      =
      \int_0^x\left(\int_0^u f(t)\,dt\right)du.
      \]

\begin{hintbox}
两边对 \(x\) 求导，初值同为 \(0\)。也可以把右边改写为三角区域上的累积。
\end{hintbox}
\end{exercisebox}

\begin{exercisebox}
\textbf{习题 23 \badge{凸函数}}

设 \(f(x)\) 在 \([0,a]\) 上二阶可导 \((a>0)\)，且 \(f''(x)\ge0\)，证明：

\[
      \int_0^a f(x)\,dx\ge af(a/2).
      \]

\begin{hintbox}
\(f''\ge0\) 表示凸性；把 \(x\) 与 \(a-x\) 配对。
\end{hintbox}
\end{exercisebox}

\begin{exercisebox}
\textbf{习题 24 \badge{凹函数与 Jensen}}

设函数 \(f(x)\) 在 \([0,1]\) 上二阶可导，且 \(f''(x)\le0\)，\(x\in[0,1]\)。证明：

\[
      \int_0^1 f(x^2)\,dx\le f(1/3).
      \]

\begin{hintbox}
凹函数满足 Jensen 反向；注意 \(\int_0^1x^2dx=1/3\)。
\end{hintbox}
\end{exercisebox}

\begin{exercisebox}
\textbf{习题 25 \badge{单调函数与振荡}}

设 \(f(x)\) 为 \([0,2\pi]\) 上的单调减少函数，证明：对任何正整数 \(n\) 成立

\[
      \int_0^{2\pi} f(x)\sin nx\,dx\ge0.
      \]

\begin{hintbox}
把区间按 \(\sin nx\) 的正负半波分组，利用单调性比较配对的函数值。
\end{hintbox}
\end{exercisebox}

\begin{exercisebox}
\textbf{习题 26 \badge{零点存在}}

设函数 \(f(x)\) 在 \([0,\pi]\) 上连续，且

\[
      \int_0^\pi f(x)\,dx=0,\qquad
      \int_0^\pi f(x)\cos x\,dx=0.
      \]

证明：在 \((0,\pi)\) 内至少存在两个不同点 \(\xi_1,\xi_2\)，使得

\[
      f(\xi_1)=f(\xi_2)=0.
      \]

\begin{hintbox}
若零点太少，\(f\) 的符号结构会受限；再结合权重 \(1\) 和 \(\cos x\) 的积分条件推出矛盾。
\end{hintbox}
\end{exercisebox}

来源：教材第七章 §3（教材页 251-267）整理；公式用 LaTeX 渲染，课后题按教材顺序录入。
